\FloatBarrier
\section{The Tested Data-Generating Process Pipeline}

A dataset is not a transcript of the world; it is the record an institution chose to keep. Standard workflows read the rows at face value and so assume the recording process is innocuous, that what was written down is a fair sample of what happened. For compliance-gated data that assumption fails: the institution layer of the structural model (Section 2.1) decides what is observed, and it does so on a variable that also drives the outcome (Section 2.2). Before modelling how the world produces outcomes, the analyst must therefore characterise how the institution decided what to record. The most direct evidence about that recording layer is the pattern of what is absent.

The TDGP pipeline operationalises this discipline. From any dataset satisfying Definition 2.1 it returns a DAG that has survived a falsification test, in three stages: read the recording layer from the data (Stage 1); hypothesise the world layer that generated the recorded values (Stage 2); and test that hypothesis against the primary documentation that defines the institution (Stage 3). The same discipline governs the analysis that the tested DAG later licenses in Sections 4 onward: the observed data is always treated as a filtered view, never as the population itself. The pipeline is dataset-agnostic; Section 4 carries it out on the Chicago Energy Benchmarking data.


\subsection{Stage 1: Shadow Matrix Analysis}

The analyst does not arrive at the data empty-handed. The documentation already supplies a hypothesis: that the dataset is compliance-gated, with reporting switched by some selection rule. What the documentation cannot supply is whether that hypothesis holds in the data at hand, and how the abstract gate maps onto the specific columns of the file. Stage 1 settles both, and it does so from the data alone, so that the recording structure is established independently of the documentary claims it will be tested against in Stage 3.

What does a gate look like in the data? A single institutional switch leaves a particular trace. Its mark is not that individual fields are missing but that they are missing together: every field the gate governs is observed for the same units and absent for the same units, because one decision controls them all. Detecting the gate is therefore a question of co-missingness: which variables vanish in concert. The natural summary of co-missingness across a dataset is the correlation between columns' missingness indicators, so that correlation, not the raw rate of missingness, is the quantity Stage 1 computes.

Concretely, for each unit $i$ and variable $j$, define the missingness indicator $R_{ij} = \mathbf{1}\{X_{ij} \text{ is missing}\}$, equal to 1 where the value is absent and 0 where it is observed. The shadow matrix is the $p \times p$ matrix of pairwise correlations among the indicator columns $R_1, \ldots, R_p$: its $(j, k)$ entry is the correlation between the missingness of variables $j$ and $k$ (the phi coefficient, for binary indicators), large when the two fields tend to be absent on the same units. Reading this matrix for structure is the whole of Stage 1.

The compliance-gate hypothesis makes sharp predictions about these correlations, and three regimes are distinguishable.

\textbf{Idiosyncratic missingness.} Off-diagonal correlations near zero mean variables go absent independently of one another. This is the signature of data missing completely at random (MCAR), under which complete-case analysis and standard imputation are unbiased (\cit{littlerubin2002}{Little \& Rubin, 2002}). A compliance-gated dataset does not look like this; observing it would refute the hypothesis.

\textbf{Block structure.} A set of variables that goes missing in lockstep, each absent almost exactly when the others are, is the empirical fingerprint of an atomic reporting gate (Condition 3 of Definition 2.1): a single switch that withholds the whole block at once. The structural model predicts this directly: because $M^{\text{obs}} = C \cdot M$ and $Y^{\text{obs}} = C \cdot Y$ are gated by the same $C$, every one of their missingness indicators reduces to $1 - C$, so the gated fields should share a single missingness pattern and correlate almost perfectly.

\textbf{Gate alignment.} Block structure shows a co-missing block; the observed $C$ shows that the compliance gate is what produces it. In the DAG, $C$ is a direct cause of the missingness ($C \to M^{\text{obs}}$, $C \to Y^{\text{obs}}$): by the institution equations $M^{\text{obs}} = C \cdot M$ and $Y^{\text{obs}} = C \cdot Y$, a field is withheld precisely when $C = 0$. The check is to correlate the observed values of $C$ against each field's missingness indicator $R_j$. Note what is being correlated: the value of $C$ against $R_j$, not the missingness of $C$ against the other indicators, since $C$ is observed for every unit and is never itself absent. For the gated fields the relationship is exact, $R_{ij} = 1 - C_i$, a correlation of $-1$; the fields whose missingness $C$ explains this way are the gated outcome and mediator set $\{M^{\text{obs}}, Y^{\text{obs}}\}$. ($C$ is observed throughout, which is also what makes the later IPW correction possible.)

Identifying the gated block does not exhaust the shadow matrix: the blanks in a real file come from several mechanisms, told apart by how they relate to $C$.

\begin{itemize}
  \item \textbf{Gate-missingness.} A field absent exactly when $C = 0$ ($R_{ij} = 1 - C_i$). This is the MAR mechanism the pipeline corrects.
  \item \textbf{Selection-side discrepancy.} Compliance does not follow the threshold cleanly: eligible units fail to report, exempt units report anyway, enforcement is uneven. Whether the observed $C$ departs from the documented rule $\mathbf{1}\{\mathbf{X}_s \in \mathcal{R}\}$ is a question for Stage 3, once that rule is in hand.
  \item \textbf{Reporting-side discrepancy.} A gated unit files partially or with error: partial submissions, measurement error, data-entry mistakes. The block's fields stop moving in lockstep, their mutual correlation dropping below one.
  \item \textbf{Structural zeros.} A quantity that does not exist for the unit, tracking its physical configuration rather than $C$. These are genuine zeros and must not be imputed.
  \item \textbf{Incidental gaps.} Random failures that track nothing. They are MCAR and handled by ordinary means.
\end{itemize}

The first three concern the gate; the realised missingness is MAR only to the extent the first dominates the next two. If the discrepancies are small, the MAR approximation and the later IPW correction hold; if large, compliance is driven by more than the threshold, and unless those extra drivers are themselves observed, selection stops being ignorable given the recorded features and IPW cannot recover the population. That selection is ignorable given the observed features, the assumption the later correction rests on (Section 3.4), is therefore not guaranteed by the documented gate but bounded by the size of this discrepancy. Stage 1 sizes each component but does not diagnose it: it cannot say which obstacle produced a discrepancy, nor whether the unexplained part depends on the withheld outcomes.

Stage 1 therefore reaches a falsifiable conclusion before any causal model is built: a block-and-gate structure is incompatible with MCAR, since under MCAR the indicators would be mutually uncorrelated (\cit{littlerubin2002}{Little \& Rubin, 2002}). Refuting MCAR is what rules out the default workflows, complete-case analysis and unconditional imputation, and forces the structural treatment that follows. It does not by itself explain why the data are absent; recasting missingness from a statistical nuisance into a causal question is the work of Stages 2 and 3 (\cit{pearl2019}{Pearl, 2019}).

The output is a map of the dataset's observation structure: the gate $C$, the gated block of fields $\{M^{\text{obs}}, Y^{\text{obs}}\}$, the structural zeros to be preserved, and the always-observed features. The map stops at this coarse partition; it does not separate the outcome from the mediators within the gated block, nor the selection variables from the other upstream features, since missingness alone cannot tell these apart. This map is the input to Stage 2, and a prerequisite for it: a causal graph over the world layer cannot be drawn until the gated fields have been identified and the structural zeros separated from withheld values, since the two enter the graph differently.

\subsection{Stage 2: DAG Hypothesis Formulation}

Stage 1 returned a map of the recording layer: the gate $C$, the co-missing block it aligns with, the structural zeros, and the always-observed features. Stage 2 turns that map into a candidate causal graph of the world layer. At this point the analyst holds only three things: the shadow matrix from Stage 1, the descriptions of the variables, and Definition 2.1. The primary documents stay closed on purpose, because they are the independent standard against which Stage 3 will try to refute the graph, and a graph built from them could not then be tested by them. The graph is nonetheless not invented from nothing: Definition 2.1 has already fixed it for the whole class, so building it for one dataset is an application of that theory rather than a fresh act of modelling.

The object being built is a directed acyclic graph (DAG): a set of nodes joined by directed edges in which no path leads from a node back to itself (\cit{pearl2009}{Pearl, 2009}; \cit{pearletal2016}{Pearl, Glymour \& Jewell, 2016}). A node is a variable, and a directed edge $A \to B$ records that $A$ is a direct cause of $B$, so that a change in $A$ produces a change in $B$ and not the reverse. Direction is what separates the graph from a symmetric table of correlations and lets it speak of intervening on a variable rather than only observing it; acyclicity orders the variables so that each is generated from those before it, which is what makes the graph a picture of a data-generating process. The DAG is the graphical face of the structural causal model of Section 2.1, and committing to it first and deriving the statistical model from it, rather than fitting a model to sample associations at the outset, is what separates this approach from purely statistical modelling (\cit{pearl2019}{Pearl, 2019}).

Building the graph under this restriction is an act of reasoning, not of look-up. The analyst sorts the variables into the roles of Definition 2.1 and argues the connections between them from what the variables mean, turning the information in hand into a structured, refutable claim about the data-generating process. This is how a DAG is built in any observational study; what differs here is only that Stage 3 will have a document to test it against. The procedure has five steps, and only the first draws on the data, even then only to a coarse partition; the rest are reasoned.

\begin{itemize}
  \item \textbf{Step 1. Carry over the recording-layer partition (data).} The gate $C$ is already fixed, the reporting status settled before the shadow matrix was built. Read against $C$, the matrix sorts the columns into the gated block, the reported fields that go missing together and in alignment with $C$, and the upstream set, the fields whose missingness is idiosyncratic or absent. This is as far as the data reaches. The matrix shows the gated block as a single undifferentiated group, with nothing to tell the outcome within it from its mediators, and it leaves the upstream set undivided in the same way.
  \item \textbf{Step 2. Name the factor of interest $Y$ (choice).} Within the gated block, the quantity the study sets out to explain or forecast is the outcome $Y$. Which field that is comes from the research question, not the data.
  \item \textbf{Step 3. Reason the mediators $M$ (conjecture).} Since the matrix cannot separate $Y$ from the other gated fields, the separation is made by argument: the remaining fields, which the descriptions present as intermediate quantities through which $Y$ is produced, are taken as the mediators feeding $Y$, giving the structural claim $\mathrm{pa}(Y) = M$.
  \item \textbf{Step 4. Reason the selection variables $\mathbf{X}_s$ (conjecture).} The upstream set is undivided too. Within it, the feature that the descriptions and the template mark as a plausible eligibility criterion, a measure of size or scale of the kind a threshold gates on (Condition 1), is taken as $\mathbf{X}_s$; the rest are $X_o$, upstream causes of the outcome that play no part in selection. Where the data records the non-compliant units, this can be corroborated by a threshold signature: $C$ turns from 0 to 1 at a cutoff in the selection variable, where a merely correlated feature would show a gradient. The signature is suggestive, not decisive, and it speaks only to selection, not to whether $\mathbf{X}_s$ causes the outcome.
  \item \textbf{Step 5. Instantiate the edges (template).} With the roles assigned, Definition 2.1 supplies every edge: $\mathbf{X}_s \to C$ and $\mathbf{X}_s \to M$, $X_o \to M$, $M \to Y$, and $M \to M^{\text{obs}} \leftarrow C$ with $Y \to Y^{\text{obs}} \leftarrow C$. No arrow is argued on its own; the template fixes them all.
\end{itemize}

The data settles only the coarse partition of Step 1, the gated block set apart from the upstream set, and at most the suggestive threshold check of Step 4. Every other commitment is categorisation and argument: which gated field is the outcome and which its mediators, and the causal claims that give the graph its content. The result is a fully specified DAG whose recording layer is grounded in the data and whose causal world layer is a reasoned conjecture, confirmed nowhere yet. The analyst cannot show, from anything now in hand, that $\mathbf{X}_s$ causes the outcome or that the mediators exhaust its causes. That is by design, and it is what keeps the construction independent of the test to come.

The manifested DAG departs from the seven-node schematic of Figure 2.1 in two ways. The first is scale: each role is a position that one variable or many may fill, so a real graph is usually larger, with several selection variables, many mediators, perhaps structure among the mediators, and more than one factor computable from $M$. The analyst keeps only the outcome the question concerns, unless several factors of interest, or a whole-network view, call for a fan $Y_1, Y_2, \ldots$, each reached by its own path $M \to Y_k$. These are expansions of the template, not departures from it: every added node still occupies one of the five roles. The second way is error. Because Steps 2 to 4 are conjectures, the assignment can be wrong: a feature placed in $X_o$ may in fact drive selection, or a field read as a mediator may be a second outcome. The graph's sharpest and most exposed claim is $\mathrm{pa}(Y) = M$, that the mediators are the only direct causes of $Y$. That claim, and the role assignments behind it, are what Stage 3 tests against the documents (Section 3.3); a divergence it detects sends the analyst back to the step that produced it.

Even before that test, the constructed graph already says what a model of $Y$ may use. The target is the latent $Y$, not its recorded form $Y^{\text{obs}}$, which exists only where $C = 1$ and so serves merely as the training signal on compliant units. The admissible inputs are the upstream features $\mathbf{X}_s$ and $X_o$, the causes above the gate that are recorded for every unit. The mediators are inadmissible twice over: because $Y = f_Y(M)$, a model given $M$ reconstructs the generating relation rather than explaining $Y$ through what a unit is; and because $M^{\text{obs}}$ is gated, it is absent for the very non-reporting units the model must reach. That near-determination is double-edged: it disqualifies the mediators as predictors of the population relationship, which Section 5 estimates from the upstream features alone, yet it is exactly what the documented formula encodes as $\mathrm{pa}(Y) = M$, the claim Stage 3 settles against the standard and the reconstruction of $Y$ from the mediators corroborates in the data (Sections 3.3, 4).

\subsection{Stage 3: DGP Falsification and Refinement}

One might expect the candidate DAG to be checked against the data, its implied independences tested in the sample. For compliance-gated data the test lies elsewhere. The structure of the graph is fixed by the documents that govern the programme, and Stage 3 opens the documents Stage 2 held back and reads it off. What the stage performs is a comparison of information: it sets each conjecture from Stage 2 beside what the documentation states and asks whether the documentation contradicts it. No statistical test of the structure is needed, because the rules that generated the data are written down.

This rests on the constitutive property of the documentation established in Section 2.3: reading the documents yields the structure itself rather than a secondary account of it. Three of their provisions settle the three conjectures Stage 2 left open.

\begin{itemize}
  \item The \textbf{selection rule} gives the actual gate: the variables it reads are $\mathbf{X}_s$, and its threshold defines $\mathcal{R}$. This bears on the Step 4 conjecture.
  \item The \textbf{reporting schema} gives the classification of the gated block: its list of required fields says which are the mediators $M$ and which the reported outcome $Y$. This bears on the Step 2 and Step 3 conjectures.
  \item The \textbf{technical standard} gives the connection from mediators to outcome: the accounting formula, emission factor, or measurement rule is $f_Y$, and its dependence on $M$ alone is the claim $\mathrm{pa}(Y) = M$.
\end{itemize}

Each conjecture is then set beside its documented counterpart. Where they agree, the assignment stands unrefuted. Where they disagree, the conjecture is refuted, and the same provision that refutes it supplies the right value in its place, so the repair is immediate: adopt the documented assignment and propagate it through the graph. A misread selection variable, a field taken for a mediator that the schema lists as a second outcome, a cause of $Y$ the formula names but the conjecture omitted, each is both caught and corrected by the document that catches it. Because the documentation entails one structure, the loop ends there, at the DAG the documents imply.

The sharpest claim, that the mediators exhaust the causes of the outcome, is settled in just this way rather than inferred from data. The reported outcome is not a separate quantity the mediators happen to track; it is computed from them by the documented formula. That $Y$ has no parent outside $M$ is thus a feature of how the institution defines the reported outcome, read from the standard that defines it, not a regularity found in the sample.

What the documents settle, they settle outright: the gate, the classification, and the outcome's parents $\mathrm{pa}(Y) = M$, the backbone of the graph. The one structural feature they do not constitute is the upstream-to-mediator edges $\mathbf{X}_s, X_o \to M$, which the correction does not rest on and the data corroborate by association (Section 2.3). One of the documented links does leave an empirical mark besides. Because the formula computes $Y$ from $M$, the reported mediators should reconstruct the reported outcome on the compliant subset; a clean reconstruction shows the recorded data following the documented formula, and a poor one shows them departing from it. This mark reaches only the $M \to Y$ link, and only where the formula is near-exact; the gate and the classification carry no comparable signal. An application runs the reconstruction as a corroborating check (Section 4), but $\mathrm{pa}(Y) = M$ is settled by the formula whether or not the check is run.

Settling the structure is not the same as guaranteeing the data obey it. What the documents fix, and fix outright, is the rule as written; what they cannot reach is the degree to which the rule was actually applied, the realized fidelity that the disturbance terms of the SCM carry and a perfectly applied procedure would not. The reconstruction residual is the outcome-side fidelity $U_Y$: a departure from the documented formula, whether from reporting error or from a contributor the documentation omits. Gate-side fidelity is the gap between observed compliance and the documented rule, legible only once that rule is read; the size of this gap is what decides whether selection stays ignorable given the observed features, the assumption the correction rests on (Section 3.4). A distortion uniform enough to imitate the rule, or one on a link with no empirical mark, escapes both the documents and the data.

The output is a tested DAG, and it is worth stating exactly what the test reached and what it did not. Three commitments were set against the written rules and left unrefuted: the selection variable and the form of the gate, the classification of the gated block into mediators and outcome, and the outcome's parents $\mathrm{pa}(Y) = M$. These are the backbone the analysis leans on, settled by an external source rather than asserted. Two further commitments the documents do not reach. The upstream-to-mediator edges $\mathbf{X}_s, X_o \to M$ are corroborated by association rather than documented, and their direction rests on the modularity the SCM assumes, as every observational graph's must. And no document certifies that realized compliance follows the documented rule closely enough for selection to be ignorable given the observed features; that ignorability is the exogeneity premise any causal model makes, assumed here as elsewhere, though here partially auditable through the gate-fidelity and manipulation checks the pipeline affords (Sections 3.1, 5). What sets the class apart is not that these residual assumptions disappear, but that they are the irreducible minimum any causal reading of the data must make: an asserted DAG assumes the backbone as well, whereas here the backbone is tested and only the standard exogeneity and extrapolation premises are left to assume. With the testable structure resting on the documentation rather than on the analyst, the full Pearlian analysis follows as a logical consequence (Section 3.4); Section 4 carries the three stages out on the Chicago Energy Benchmarking data.

\subsection{What the Tested DAG Licenses}

It is worth naming the question the tested DAG is meant to answer before reaching for the tools that answer it. We want to learn about the outcome $Y$ in two senses: which of the available factors explain its variation, and which predict it. The material to learn from is the observable set $\mathcal{O} = \{\mathbf{X}_s, X_o, C, M^{\text{obs}}, Y^{\text{obs}}\}$, and in it the outcome and the mediators appear only in their gated forms, $Y^{\text{obs}} = C \cdot Y$ and $M^{\text{obs}} = C \cdot M$, present for compliant units and absent for the rest.

An analysis that does not keep the DAG in view reads $Y^{\text{obs}}$ and $M^{\text{obs}}$ as though they were $Y$ and $M$. It then answers both questions for the compliant subsample while taking the answers for the population's, which is the error the whole pipeline exists to prevent. Holding the DAG in view forbids the substitution: $Y^{\text{obs}}$ is a gated shadow of the factor of interest, not the factor, so the task is to learn about a latent quantity from its observable traces. Pearl's ladder of causation is the means. Its three rungs, seeing, doing, and imagining, are each a stronger question than the one below, and each was closed while the graph was only asserted (\cit{pearl2009}{Pearl, 2009}). The climb begins with reading the graph.

At the lowest rung a DAG is a map of association, and d-separation is the rule for reading it (\cit{pearl2009}{Pearl, 2009}; \cit{pearletal2016}{Pearl, Glymour \& Jewell, 2016, §2.4}): information travels a path unless the path is blocked, a chain or a fork is blocked by conditioning on its middle node, and a collider is blocked until its middle node, or a descendant of it, is conditioned on. The taxonomy is textbook material and is not rehearsed here; what matters is which of its patterns the gate creates. $\mathbf{X}_s \to M \to Y$ is a chain, cut by conditioning on the mediators, $\mathbf{X}_s \perp Y \mid M$. $C \leftarrow \mathbf{X}_s \to M$ is a fork, cut by conditioning on the selection variables, $C \perp M \mid \mathbf{X}_s$. And the analyst's records are colliders, $M \to M^{\text{obs}} \leftarrow C$, which conditioning opens rather than cuts. These three readings are the entire associational apparatus the argument below uses.

The second rung asks not what accompanies what, but what follows from acting. The intervention $do(X = x)$ sets $X$ from outside the system, severing the arrows into it so its own causes no longer bear on it, and the distribution it produces, $P(Y \mid do(x))$, is in general not the observed $P(Y \mid x)$. The back-door criterion certifies when the one is recoverable from the other: a set $Z$ that blocks every path from $X$ to $Y$ running against the arrows, and contains no effect of $X$, recovers the effect by adjustment, $P(Y \mid do(x)) = \sum_z P(Y \mid x, z)\, P(z)$ (\cit{pearletal2016}{Pearl, Glymour \& Jewell, 2016, §3.3}).

Inverse-probability weighting is the back-door adjustment re-expressed, and the re-expression is what licenses it. Multiply and divide each term of the adjustment formula by the propensity $P(x \mid z)$, and recognise the numerator that results as the observed joint distribution:
\[
P(Y \mid do(x)) = \sum_z P(Y \mid x, z)\, P(z) = \sum_z \frac{P(Y, x, z)}{P(x \mid z)}.
\]
The same quantity is therefore obtained by counting each observed case $(Y, x, z)$ with weight $1 / P(x \mid z)$ (\cit{pearletal2016}{Pearl, Glymour \& Jewell, 2016, §3.6}): reweighting a sample this way makes it behave as though drawn under the intervention. Because the two expressions are equal, the weighting inherits the back-door criterion exactly. It recovers the effect when $Z$ closes the back door, and injects bias of its own when $Z$ does not (\cit{pearletal2016}{Pearl, Glymour \& Jewell, 2016, p. 75}).

This paper turns the device on the gate, with $C$ now playing the treatment and the upstream features the adjusting set $Z$. The two upstream features enter that set for different reasons, and only one of them does causal work.

\textbf{The role of $\mathbf{X}_s$: the confounder, and the source of the bias.}

\begin{itemize}
  \item It occupies two positions at once: it sets compliance through the gate ($\mathbf{X}_s \to C$) and causes the outcome through the path $\mathbf{X}_s \to M \to Y$. A common cause of selection and outcome confounds their association, which is why the compliant subsample is enriched on $Y$.
  \item This dual position is Condition 2 doing its work. Were $\mathbf{X}_s$ not a cause of $M$, there would be no enrichment and nothing to correct: the edge into $M$ is what makes the bias exist.
  \item Conditioning on $\mathbf{X}_s$ removes the confounding. In the graph, every route from $C$ to $Y$ is the fork $C \leftarrow \mathbf{X}_s \to M \to Y$ or a shut collider at a recorded field, and conditioning on $\mathbf{X}_s$ closes the fork while leaving the colliders shut, giving $C \perp Y \mid \mathbf{X}_s$. The two facts this uses, that the gate reads $\mathbf{X}_s$ and that $\mathrm{pa}(Y) = M$ leaves the outcome no other parent by which a path could reopen, are the ones Stage 3 settled against the documents.
\end{itemize}

\textbf{The role of $X_o$: not a confounder, a safe default.}

\begin{itemize}
  \item Nothing connects $X_o$ to $C$, so it can never confound the $C$-$Y$ association, and the independence above holds with or without it.
  \item It is conditioned on for two milder reasons: it causes the outcome, so the analysis wants it in the population relationship $P(Y \mid \mathbf{X}_s, X_o)$ it is after, and a pre-treatment feature the gate ignores can neither open a path nor distort the adjustment.
  \item Its inclusion is therefore safe in either case: if $X_o \to M$ holds it sharpens the estimate, and if it does not the feature is simply inert.
\end{itemize}

The blocking just described holds within the graph, and it carries one assumption the documents do not discharge. It supposes the graph is complete on the selection side, that no unobserved driver of compliance is correlated with the outcome once the upstream features are fixed. Under the documented gate, compliance is a function of $\mathbf{X}_s$ alone and this holds by construction. Real programmes comply imperfectly, so wherever realized compliance is not fully determined by $\mathbf{X}_s$, the condition becomes the selection-ignorability assumption standard to every inverse-probability correction (\cit{hernanrobins2020}{Hernán \& Robins, 2020}): not settled by the documents, which fix only which variable selection keys on, but partially auditable here through the gate-fidelity and manipulation diagnostics the pipeline affords (Sections 3.1, 5). Naming it is what keeps the tested structure honest, for it is the residual the reading of documents cannot reach.

With the back door closed under this assumption, weighting each compliant unit $i$ by the inverse of its selection propensity, $w_i = 1 / P(C = 1 \mid \mathbf{X}_{s,i}, X_{o,i})$, restores the missing mass and lets the compliant sample stand for the full population (\cit{rosenbaumrubin1983}{Rosenbaum \& Rubin, 1983}; \cit{bareinboimpearl2012}{Bareinboim \& Pearl, 2012}). The corroborated edges $\mathbf{X}_s, X_o \to M$ enter none of the weighting: they make the bias exist, not the correction valid, so should the upstream features shape $M$ more weakly than supposed, there is less bias for the weights to remove and the correction stays sound rather than turning wrong. Section 5 estimates the propensity and applies the weights.

This recovery stops at a boundary the weighting cannot cross. Reweighting recovers the population only where the gate leaves overlap, where compliant units exist to be up-weighted; under a deterministic gate there are none below the threshold, $P(C = 1 \mid \mathbf{X}_s) = 0$, and no weight stands in for units the data never recorded. What reaches that region is not the weighting but the structural form, extended by the modularity of its equations: a mechanism holds whatever the distribution of its inputs, so an identified $f_M$ and $f_Y$ carry past the cutoff. The tested structure narrows what this asks one to believe. The $M \to Y$ link is the documented $f_Y$, invariant by construction, so it extrapolates for free; only the learned $f_M$, the upstream-to-mediator step $\mathbf{X}_s, X_o \to M$, rests on the assumption that its form is stable below the threshold, an assumption checkable at least at the edge of the observed range. The weighting therefore de-biases within the support, and the tested structure extends the reach beyond it, which is where a graph that has survived a test earns more than an asserted one: both must extrapolate, but only the tested one does so with a warrant.

The third rung asks not for new data but for the model the second rung delivered. Intervention is what makes the population's mechanism knowable: only once selection is corrected can the structural equations of Definition 2.1 be estimated without bias, $f_M$ carrying the upstream features into the mediators and $f_Y$ the mediators into the outcome. A counterfactual takes that known model and runs it again for one unit, under a value the unit did not have, to ask what its outcome would then have been. The answer comes entirely from the equations, which is why it waits on a structure that has survived Stage 3.

Take a compliant unit, with $\mathbf{X}_s, X_o, M, Y$ all observed, and ask what its outcome would have been had its non-selection features $X_o$ taken a value $x$ instead. The computation has three steps (\cit{pearl2009}{Pearl, 2009}; \cit{pearletal2016}{Pearl, Glymour \& Jewell, 2016, Ch. 4}).

\begin{itemize}
  \item \textbf{Abduction.} Recover the unit's own disturbances from what it showed, the $u_M$ that solves $M = f_M(\mathbf{X}_s, X_o, u_M)$ and the $u_Y$ that solves $Y = f_Y(M, u_Y)$; in an estimated model these come as a posterior, $P(u \mid \text{evidence})$.
  \item \textbf{Action.} Replace $X_o$ with the counterfactual value $x$, holding $\mathbf{X}_s$, the disturbances, and the equations fixed. Because $X_o$ does not enter the gate $C = \mathbf{1}\{\mathbf{X}_s \in \mathcal{R}\}$, the unit's compliance is untouched, and the query concerns the world layer alone.
  \item \textbf{Prediction.} Re-run the model, $Y_{x} = f_Y\big(f_M(\mathbf{X}_s,\, x,\, u_M),\; u_Y\big)$, carrying $x$ through the mediators to the outcome.
\end{itemize}

The disturbances held over from abduction are what make the answer the unit's own rather than the population's average, and the equations held over from intervention are what make the prediction computable at all. Section 7 evaluates two such queries.

The apparatus is procedural and epistemological, not statistical. The ladder fixes what must be done, which variables to condition on, which to leave alone, how to weight, but it names no estimator and assumes no distributional form; it is theory that makes room for statistics rather than statistics itself. A frequentist regression, a Bayesian hierarchical model, and a machine-learning ensemble are all admissible ways to study $Y$, and each removes the selection bias so long as it obeys the rules the arguments lay down: target the latent outcome and not its gated shadow, take inputs from the upstream features alone, and carry the inverse-probability weights. Which of them explains and predicts best, once the causal ground is shared, is a separate question.

Section 4 carries the three stages out on the Chicago Energy Benchmarking data, and Sections 5 to 7 take up that separate question across the model families, climbing the three rungs in turn.

